% October 9, 2026 build (agent draft): model appendices following the deck's
% appendix (supply arms, CES variant, solution algorithms). Tables for the CES
% variant are the deck's generated files (latex/JMP_slides/ces_tables/).
\section{Housing Supply Arms}
\label{app:supply}

Figure~\ref{fig:app-slow-transition} compares the estimated transition of
Section~\ref{sec:quant-transition}, in which the stock stays on the supply
curve, with the same path when the stock adjusts slowly from 2023, as in
\eqref{eq:policy-slow}. With the slow stock, households fall less in the
first century, total housing falls more slowly, and the price falls faster at
first; the two paths reach the same stationary equilibrium.
Figures~\ref{fig:app-slow-policies} and~\ref{fig:app-frozen-policies} show the
paths of the policies of Section~\ref{sec:policy} with the slowly adjusting
and the frozen stock.
\mocktodo{the deck's own note: refit the 2007--2023 transition under the
slowly adjusting stock for these figures (the re-estimated pair of
Table~\ref{tab:policy-channels} exists for the tax arms only).}

\begin{figure}[p]
\centering
\includegraphics[width=\linewidth,trim=0 0 0 50,clip]{../../output/model/transition_h100_14402_20261007/appendix_supply_arms_20261009/slow_stock_vs_estimated.pdf}
\caption{Estimated transition and a slowly adjusting stock from 2023}
\label{fig:app-slow-transition}
\end{figure}

\begin{figure}[p]
\centering
\includegraphics[width=\linewidth,trim=0 0 0 50,clip]{../../output/model/transition_h100_14402_20261007/appendix_supply_arms_20261009/slow_stock_policies.pdf}
\caption{Policies from 2023 with a slowly adjusting stock}
\label{fig:app-slow-policies}
\par\medskip
\includegraphics[width=\linewidth,trim=0 0 0 50,clip]{../../output/model/transition_h100_14402_20261007/appendix_supply_arms_20261009/frozen_stock_policies.pdf}
\caption{Policies from 2023 with the stock frozen at its 2023 level}
\label{fig:app-frozen-policies}
\end{figure}

\section{A CES Housing Composite}
\label{app:ces}

The baseline combines a Cobb--Douglas composite with the space requirement
$h_P$. As a robustness check I replace both with a CES composite in which
children raise the rooms needed for given housing services:
\begin{equation}
 u_t=\frac{1}{1-\sigma}\left(\frac{Q_t}{e(m_t)}\right)^{1-\sigma}+\psi_t m_t^{1-\gamma},\qquad
 Q_t=\Bigl[\alpha_0\,c_t^{\rho}+(1-\alpha_0)\Bigl(\frac{\chi^{\1\{\mathrm{own}\}}\,h_t}{k(m_t)}\Bigr)^{\rho}\Bigr]^{1/\rho},
\end{equation}
with $\rho=(\varepsilon-1)/\varepsilon$, $\varepsilon=0.67$, and
$k(m_t)=\exp\bigl[\eta_1\1\{m_t\ge1\}+\eta_2\max(m_t-1,0)\bigr]$. Housing and
consumption are complements, and without a floor space is not a fixed entry
cost of a child. I estimate $\eta_1$ and $\eta_2$, the rooms needed by the
first and each later child at home, and $\alpha_0$, so twelve parameters are
chosen to match thirteen moments. The new moments are the rent share of
renters aged 30--55 with no child at home (0.360, CE), the probability of
moving from renting to owning within four years at ages 22--41 (0.281, PSID),
and the share of mothers with three or more children (0.357, CPS). The rooms
added at the first birth are re-measured (1.025, PSID), the room gap between
larger and smaller families becomes a target, and the recent-parent
ownership gap becomes a validation moment.

% completed fertility at 46--49, CES at the calibrated point, fixed price and rebate (generated)
\def\cesLTVhundred{$+0.1\%$}
\def\cesLowerRisk{$+13.2\%$}
\def\cesCapTen{$+0.5\%$}

Table~\ref{tab:app-ces-fit} reports the fit. It is worse than the
benchmark's: slightly better on childlessness, family size, rooms and the
rooms at first birth, and worse on timing, ownership, the rent share, the
family room gap, wealth and bequests. The responses of completed fertility
are as in the benchmark: \cesLTVhundred{} with 100 percent financing,
\cesLowerRisk{} with lower earnings risk, and \cesCapTen{} with a ten-room
rental cap.

\begin{table}[t]
\centering
\caption{CES variant: targets and fit}
\label{tab:app-ces-fit}
\small
\begin{minipage}[t]{0.66\linewidth}\vspace{0pt}
\begin{tabularx}{\textwidth}{@{}Xrrr@{}}
\toprule
Moment & Target & CES & Cobb--Douglas \\
\midrule
Childless, women 40--44 (\%) & 19.8 & 20.1 & 20.2 \\
One child among mothers, 40--44 (\%) & 21.4 & 21.5 & 21.7 \\
Mothers with 3+ children, 40--44 (\%) & 35.7 & 38.7 & 39.5 \\
Mean age at first birth (years) & 25.98 & 26.02 & 25.96 \\
Children ever born by age 25 & 0.810 & 0.542 & 0.553 \\
\addlinespace[3pt]
Mean occupied rooms & 5.73 & 5.80 & 5.82 \\
Ownership, heads 30--55 (\%) & 67.6 & 66.6 & 67.0 \\
Renter to owner within 4 years (\%) & 28.1 & 34.7 & 31.0 \\
Rent share, childless renters (\%) & 36.0 & 17.2 & 23.3 \\
First-birth room response & 1.025 & 0.801 & 1.264 \\
Room gap, 3+ vs 1--2 children at home & 0.385 & 0.538 & 0.294 \\
\addlinespace[3pt]
Wealth / annual gross earnings & 4.46 & 5.07 & 4.58 \\
Annual bequests / wealth (\%) & 0.73 & 0.86 & 0.71 \\
\bottomrule
\end{tabularx}

\end{minipage}\hfill
\begin{minipage}[t]{0.3\linewidth}\vspace{0pt}
\begin{tabular}{@{}lr@{}}
\toprule
Parameter & Value \\
\midrule
$\beta_{\rm annual}$ & 0.948 \\ % Annual discount factor
$\kappa_1$ & 0.100 \\ % First-birth taste scale
$\kappa_C$ & 0.246 \\ % Later-birth taste scale
$\chi$ & 1.061 \\ % Owner housing-service multiplier
$H_0$ & 7.941 \\ % Housing supply scale
$\theta_0$ & 0.421 \\ % Bequest-motive scale
$\xi$ & 0.304 \\ % First-birth utility cost
$\psi_0$ & 0.150 \\ % Initial preference for children
$\gamma$ & 0.156 \\ % Curvature of the benefit of children
$\kappa_T$ & 0.016 \\ % Tenure taste scale
$\alpha_0$ & 0.720 \\ % Consumption weight in the composite
$\eta_1$ & 0.630 \\ % Housing need of the first child
$\eta_2$ & 0.059 \\ % Housing need of each later child
\bottomrule
\end{tabular}

\end{minipage}
\end{table}

Before recalibration, Table~\ref{tab:app-ces-credit} re-fits only $\psi$ to
19 percent childlessness under each specification, at fixed prices and the
six-room rental cap. Mortgage credit moves completed fertility by between $-2$
and $+0.3$ percent in all thirteen specifications that reach 19 percent
childlessness, and the negative sign comes mostly from the floor. Credit to
renters raises births on impact by 6 to 9 percent and completed fertility by
between $-0.6$ and $+2.4$ percent in the cases shown.

\begin{table}[t]
\centering
\caption{CES variant: credit across specifications}
\label{tab:app-ces-credit}
\small\setlength{\tabcolsep}{4pt}
\begin{tabularx}{\textwidth}{@{}Xccc rr rr@{}}
\toprule
 & & & & \multicolumn{2}{c}{Financed share $80\%\to95\%$} & \multicolumn{2}{c}{Renter credit, 2 years} \\
\cmidrule(lr){5-6}\cmidrule(l){7-8}
 & $\varepsilon$ & Floor & Need $k$ & Impact & Completed & Impact & Completed \\
\midrule
Cobb--Douglas (benchmark) & 1 & yes & 1 & $+0.42\%$ & $-1.35\%$ & $+8.5\%$ & $+0.62\%$ \\
CES & 0.5 & yes & 1 & $-0.16\%$ & $-1.40\%$ & $+8.0\%$ & $+0.22\%$ \\
CES & 0.25 & yes & 1 & $-0.06\%$ & $-0.87\%$ & $+7.3\%$ & $-0.25\%$ \\
CES & 0.5 & yes & 1.5 & $+0.05\%$ & $-1.10\%$ & $+7.0\%$ & $-0.63\%$ \\
CES, no floor & 0.5 & no & 1 & $+0.44\%$ & $-0.03\%$ & $+6.9\%$ & $+2.25\%$ \\
CES, no floor & 0.5 & no & 1.5 & $+0.51\%$ & $-0.18\%$ & $+8.3\%$ & $+2.39\%$ \\
Cobb--Douglas, no floor & 1 & no & 1.5 & $+1.03\%$ & $+0.23\%$ & $+8.3\%$ & $+2.42\%$ \\
\bottomrule
\end{tabularx}
\par\smallskip
\begin{minipage}{0.95\linewidth}
\footnotesize\textit{Notes:} Each row re-fits only $\psi$ to 19 percent
childlessness; house prices fixed, rental cap of six rooms.
\end{minipage}
\end{table}

\section{Solution Algorithms}
\label{app:algorithms}

\paragraph{Stationary equilibrium.} I solve for the house price $P$, the
pension $\varpi$ and the rebate $T$, given parameters and a fixed preference
for children.
\begin{enumerate}
\item Guess $(P,\varpi,T)$ and derive the rent $r$ from the user cost.
\item Solve the household problem backward over age and compute the choice
probabilities.
\item Construct the stationary distribution of households. Set the mass of
households equal to housing supply divided by mean housing per household.
\item Compute three residuals: excess demand for housing, the pension budget
and the rebate budget. The loss is the largest residual in absolute value.
\item Update $(P,\varpi,T)$ and repeat until the loss is below tolerance.
Then apply the population law for one period and verify that the
distribution reproduces itself.
\end{enumerate}

\paragraph{Transition.} I solve for the paths $(P_t,\varpi_t,T_t)_{t\geq t_0}$,
given the inherited households and a preference believed permanent.
\begin{enumerate}
\item Solve the new stationary equilibrium; it supplies the terminal value
function.
\item Guess the three paths and derive rents from \eqref{eq:m-rent}.
\item Solve the household problem backward in calendar time from the terminal
value.
\item Advance households and the queue of entrants forward from the inherited
state, and aggregate housing demand and taxes at each date.
\item Compute the three residuals at every date. The loss is the largest
residual across dates.
\item Update the paths jointly and repeat until the loss is below tolerance.
Extend the horizon to check that early responses are stable.
\end{enumerate}
The transitions in the paper use 100 four-year periods, with tolerances of
$2\times10^{-4}$ on the housing market and $2\times10^{-5}$ on the pension and
the rebate.
